\chapter{Roles and user journeys}\label{ch:journeys}
\chapterstatus{complete}{2026-11-17}{\src{src/app}, \src{src/server/services}}

How each role uses Slotty from start to finish, with the screen, the service that runs, and what the user gets back. Update this chapter whenever a flow changes.

Roles are \textbf{per course} (\code{Enrollment.role}). A TA can be a student in another course (their home page then has separate \ui{Teaching} and \ui{Studying} sections), but an instructor is never a TA or a student anywhere, and students are never instructors. Only admins name instructors. The instructor manages the course; a TA manages it only when it has no instructor. Instructors usually don't take demos: they give the final check on marks. \textbf{Admin} is a separate, platform-wide role: an admin is never a member of any course.

Times show in the course timezone (Pakistan time for now) on each user's 12- or 24-hour clock (account page).

\section{Visitor (not signed in)}\label{sec:journey-visitor}

Accounts are invite-only, so most people first meet Slotty through an invite email or a link from
course staff. Signed-in people opening \route{/} go straight to their home (\route{/dashboard},
which sends admins to \route{/admin} and platform accounts to \route{/platform}).

{\tablefont
\begin{xltabular}{\linewidth}{@{}l P{34mm} P{30mm} L@{}}
  \toprule
  \thead{\#} & \thead{Step and screen} & \thead{Service} & \thead{What happens / feedback}\\
  \midrule
  \endhead
  V1 & Finds out what Slotty is\newline{\color{muted}\footnotesize \route{/}} & --- (static, \src{src/components/landing}) & The landing page, styled as a departure board for demo slots. The hero shows one student booking one demo (Lab~4 of CS~2001): a pointer books 10:15 and the row flips to \ui{BOOKED}. Then: the shared-spreadsheet way demo week used to run; the five steps of one demo (a stepper that follows the scroll, with \ui{Skip}); what students, TAs and instructors each see of the same slot; the details; privacy in four facts with a link to \route{/privacy}; who makes Slotty and that it is open source; the eight questions confused users ask (four up front, four behind \ui{More questions}); and a final \ui{Sign in}. Every call to action is \ui{Sign in} or the support address: there is no sign-up. Search engines may index only \route{/} and \route{/privacy} (\src{src/app/robots.ts}).\\
  V2 & Gets help\newline{\color{muted}\footnotesize \code{support@slotty.msaeedsaeedi.com}} & SES receiving (\cref{sec:aws-email}) & The support address is on the landing page, the sign-in pages and in invite emails. Mail to it is forwarded to the product owner, with the sender as \code{Reply-To}. A new course starts here: a TA can then create it, or an admin sets it up with its instructor.\\
  \bottomrule
\end{xltabular}}

\section{Student}

{\tablefont
\begin{xltabular}{\linewidth}{@{}l P{34mm} P{30mm} L@{}}
  \toprule
  \thead{\#} & \thead{Step and screen} & \thead{Service} & \thead{What happens / feedback}\\
  \midrule
  \endhead
  S1 & Gets invited\newline{\color{muted}\footnotesize Email → \route{/invite/[token]}} & \code{accounts.acceptInvite} & Staff import the class list, which sends an invite email valid for 14 days. The student sets a name and password, and the account becomes \code{ACTIVE}. If they were already active, they get a ``Added to <course>'' notification instead. The invite email links to the landing page (V1) and the support address.\\
  S2 & Signs in\newline{\color{muted}\footnotesize \route{/login}, \route{/forgot-password}} & \code{accounts.authenticate} & A clear message for a wrong password, a disabled account, or an invite not yet accepted. Password reset links last 2 hours. Under the form, \ui{No account yet?} explains that accounts come by invite, that \ui{Forgot password} sends a fresh invite to an account not yet activated, and gives the support address; every sign-in page links to the landing page.\\
  S3 & Sees what to do\newline{\color{muted}\footnotesize \route{/dashboard} (\ui{Studying})} & \code{home.getHome} & Next demo first (time, room, TA, \ui{Join online}, \ui{Add to calendar}), then a to-do list: book a demo (urgent when the window closes within 3 days), booking opens soon, a missed demo, marks released. Then their courses.\\
  S4 & Opens a course\newline{\color{muted}\footnotesize \route{/courses/[id]}} & \code{listAssignments} & Published and closed assignments with a status: \ui{Not booked}, \ui{Booked}, \ui{Closed}, \ui{Marked x/y}.\\
  S5 & Books a slot\newline{\color{muted}\footnotesize \route{/courses/[id]/assignments/[aid]}} & \code{bookings.bookSlot} & Slots grouped by day in the course timezone. Full slots are hidden and counted. A disabled \textbf{Book} button shows why it's disabled; each button's accessible name includes the day, time, venue and host. Success: a toast, and a confirmation email a few minutes later for the final state (no in-app item for the student's own action).\\
  S6 & Gets reminders\newline{\color{muted}\footnotesize Email + bell} & \code{reminders.queueDueReminders} (worker) & 24h and 1h before the slot. Skipped if the student booked inside that window.\\
  S7 & Reschedules\newline{\color{muted}\footnotesize Same page, \code{?reschedule=1}} & \code{bookings.rescheduleBooking} & Picks a new slot and confirms. The old booking is cancelled and a new one created. Uses one change. Shows ``N of M changes left'' and the exact time changes lock. Offered only while a move is possible; the history labels the old booking ``Rescheduled''.\\
  S8 & Cancels\newline{\color{muted}\footnotesize Same page} & \code{bookings.cancelBooking} & The confirmation states the cost (``This uses your last change\dots{}''). Only allowed while booking is open, \code{allowStudentCancel} is on, the booking is outside the freeze window, and the demo has no marks. Cancelling with no changes left doesn't spend one, but the student can't rebook themselves until staff place them or grant a change; the page then shows only the reason and the request box.\\
  S9 & Handles staff changes\newline{\color{muted}\footnotesize Bell + email} & --- & \code{slot.cancelled}, \code{booking.cancelled_by_staff} (always with a reason), \code{slot.venue_changed}, \code{assignment.rules_changed}. Rebooking after a staff cancellation doesn't use a change.\\
  S10 & Attends the demo & --- & The booking card shows venue, meeting link and TA.\\
  S11 & Sees results\newline{\color{muted}\footnotesize Assignment page} & \code{evaluations.getMyResult} & Only released marks are shown: total, rubric rows, feedback. While staff correct unlocked marks, the student keeps seeing the released ones (marked as being corrected) and any marks question. Private notes are never shown. \code{evaluation.finalized} notification.\\
  S12b & Gets help when stuck\newline{\color{muted}\footnotesize Assignment page} & \code{requests.createRequest}, \code{waitlist.joinWaitlist} & ``Need a different time? Ask course staff'' opens automatically when self-service is blocked (freeze window, no changes left, closed, nothing free). With every slot full, ``Notify me when a slot frees up''. After marks are released, ``Question about your marks?''.\\
  S12c & Calendar \& account\newline{\color{muted}\footnotesize Booking card, \route{/account}} & \code{calendar.bookingCalendarOptions}, \code{accounts.*} & \ui{Add to calendar} on a booking opens Google Calendar or Outlook with the demo filled in, or downloads an \code{.ics} (Apple and others). Automatic calendar sync is shown as \ui{coming soon}. The account page shows only what applies to the role: clock (12/24 h), reminder emails, device notifications, password, devices.\\
  S12d & Installs the app / goes offline\newline{\color{muted}\footnotesize Dashboard hint, account page} & PWA (\src{src/components/pwa.tsx}, \src{public/sw.js}) & Installs Slotty to the home screen. Turns on notifications for the device (asked after the first booking, or from the account page). Offline, recently opened pages still show the demo time, room, link and TA, with a banner. Changes wait until the device reconnects.\\
  S12 & Reviews history\newline{\color{muted}\footnotesize \route{/bookings}, \route{/notifications}} & \code{listMyBookings}, \code{inbox} & Every booking, including cancelled and past ones; moved bookings say ``Rescheduled'' or ``Moved by staff''. The notification feed is grouped by assignment, the bell counts only items that need the person, and it has ``Mark all read''.\\
  \bottomrule
\end{xltabular}}

\section{TA (can run a course alone)}

{\tablefont
\begin{xltabular}{\linewidth}{@{}l P{34mm} P{30mm} L@{}}
  \toprule
  \thead{\#} & \thead{Step and screen} & \thead{Service} & \thead{What happens / feedback}\\
  \midrule
  \endhead
  T0 & Sees what to do\newline{\color{muted}\footnotesize \route{/dashboard} (\ui{Teaching})} & \code{home.getHome}, \code{getDemoDay} & Today's demos, then a to-do list across their courses: past demos missing attendance, marks to finish or returned for changes, student requests, submissions to review (instructors), assignments to add slots to or open. Each item goes straight to the place to do it.\\
  T1 & Creates a course\newline{\color{muted}\footnotesize \route{/courses/new}} & \code{courses.createCourse} & Code, title, term; timezone is Pakistan time for now. The creator is the course's TA. Open to TAs and to new staff not yet in a course; students and instructors don't see \ui{New course}. Courses with an instructor are created by an admin from the console.\\
  T2 & Adds people\newline{\color{muted}\footnotesize \code{…/manage/roster}} & \code{addMember}, \code{previewRoster} → \code{importRoster} & \ui{Add someone}: one email, name (for new people), role, section. \ui{Import a class list}: the columns are shown with a template to download; upload a CSV (Google Classroom exports work), paste cells copied from a spreadsheet, or paste one email per line. A preview shows who is new, added, changing role or already there, and which lines are skipped (e.g. an admin's email). Only admins can add instructors (the role isn't offered to anyone else). Adding someone already there with a blank section keeps their section. Students added late are told which demos are open. Removing a student releases their upcoming bookings. Removing a host is blocked while they host upcoming slots.\\
  T3 & Adds venues\newline{\color{muted}\footnotesize \code{…/manage/venues}} & \code{createVenue}, \code{deleteVenue} & Name, location, meeting link. Creating a Google Meet link automatically is shown as \ui{coming soon}. A venue can't be deleted while live slots use it.\\
  T4 & Creates an assignment\newline{\color{muted}\footnotesize \code{…/manage/assignments/new}} & \code{assignments.createAssignment} & \ui{What}, \ui{Marking} (max marks and rubric rows, which must add up to exactly the max marks; the running total turns red until they do), \ui{When} (demo window, slot length, break, students per slot), and \ui{Booking rules} folded behind a one-line summary. Every field is checked as you type with the same limits as the server; nothing typed is lost when a save fails. A side summary says what the settings mean and what comes next.\\
  T5 & Adds slots\newline{\color{muted}\footnotesize Assignment → \ui{Slots} tab, \ui{Add slots}} & \code{slots.addAvailability} & A box at the top says how many slots are still needed: students still to book against free places in upcoming slots, and whether the days picked cover them. Pick one or more days (each shows how many slots it already has in this assignment, and as \ui{busy} the host's slots in other assignments, named in a note and the button's label; an empty course says \ui{No students yet}), the hours, and a venue; a preview lists the slots each day will get and how many are skipped because the host is already busy. Slots can be added in rounds as availability is known. TAs add slots for themselves; instructors add them for any staff member (a TA by default). Slots stay hidden until booking opens; once it's open, students who haven't booked are told about new times.\\
  T6 & Opens booking\newline{\color{muted}\footnotesize \ui{Booking} panel on the assignment} & \code{publishAssignment} & The status says exactly where booking stands: \ui{Not open yet}, \ui{Opens Fri 9 Oct, 9:00 AM} (open but starting later), \ui{Booking open}, \ui{Booking closed}; students see the same on their course page. Opening needs slots first, because every student is notified. Students are notified now, or told the opening time (and told again, as a \ui{new booking time}, if it moves).\\
  T7 & Manages slots\newline{\color{muted}\footnotesize \ui{Slots} tab} & \code{changeVenue}, \code{reassignHost}, \code{deleteUnbookedSlots}, \code{updateSlotCapacity}, \code{cancelSlot}, \code{staffCancelBooking} & A table with column headers (time, host, venue, students), grouped by day; days fold away (past days start folded) and the table scrolls on its own. Selecting slots brings up a floating bar: picking a venue or host applies at once and says how many students will be told; delete asks first. Each row's menu changes students per slot or cancels the slot, in dialogs that close with Esc or a click outside. A reason is required whenever a student is affected. Only instructors can hand slots to another host.\\
  T7b & Edits the assignment\newline{\color{muted}\footnotesize \ui{Edit settings} (next to the settings summary)} & \code{updateAssignment}, \code{deleteAssignment} & A notice explains the impact on booked students. The edit is blocked if max marks drop below awarded marks or the window no longer covers booked slots. Booked students are told about rule changes. Deleting (only before anyone has booked) is at the bottom of this page.\\
  T7c & Handles student requests\newline{\color{muted}\footnotesize \code{…/manage/requests} (badge in nav)} & \code{requests.resolveRequest} & Open the student to move them, clear a no-show, or grant an exception, then mark the request handled or declined with a reply. Moving the student keeps the request open for that reply, and each card shows the student's current booking.\\
  T7d & Fixes one student's booking or marks\newline{\color{muted}\footnotesize Student page (\code{…/evaluate/[studentId]})} & \code{saveDraft}, \code{clearMarks}, \code{staffPlaceStudent}, \code{allowRebookAfterNoShow}, \code{setAllowance} & Marks can be entered right here (same editor as the marking sheet, with its own \ui{Submit for review} / \ui{Release marks}). \ui{Clear marks} removes draft marks, with a reason. The \ui{Booking} card places or moves them into any upcoming slot (the confirmation names the new slot), even inside the freeze window, and optionally over capacity, unless marks are already recorded. None of it uses the student's changes.\\
  T8 & Runs the day\newline{\color{muted}\footnotesize \route{/today} (all courses; \ui{Demo day} in the header), or a course's \ui{Today} tab} & \code{demo-day.getDemoDay}, \code{markAttendance} & Opens on who's \textbf{now} or \textbf{up next}. The day's demos are grouped by assignment, each group linking to its course and assignment. Every booked demo has a coloured status in words (\ui{Booked}, \ui{In progress}, \ui{Attendance missing}, \ui{To mark}, \ui{Marked}, \ui{No-show}); open slots are one quiet line. The week strip colours each day by its most pressing state, with a legend (attendance missing, marks to finish, demos booked, all done) and pages a week at a time. \ui{Present → mark} records attendance and opens the marking sheet. \ui{Still to finish} lists earlier days without attendance and, per assignment, demos still to mark. Refreshes itself every 30 s. ``My demos'' by default; ``Everyone'' shows all hosts.\\
  T9 & Marks\newline{\color{muted}\footnotesize Marking sheet (\code{…/assignments/[aid]/mark}) or the student page} & \code{getMarkingSheet}, \code{saveDraft}, \code{sheetAttendanceAction} & One student at a time with \ui{Back}/\ui{Next} (also Alt+←/→, Ctrl+Enter), and the day's demos as a compact sheet underneath; clicking a row opens that student. A thin segmented bar shows every demo's state and fills the current one as its time runs. Rubric scores sit in one row of cells (Enter moves to the next), total is live (shown as ``--'' while any value is invalid), out-of-range values are flagged and not saved. Everything autosaves as a draft. Entering a mark records the student as present; \ui{No-show} asks first. \ui{Follow the clock} moves to the next demo when time is up. Marking before the demo has started, or a student who never booked, asks for a reason first, which is kept and audited.\\
  T10 & Submits\newline{\color{muted}\footnotesize Marking sheet (once the day's demos are over) or \ui{Students \& marks} tab} & \code{submitEvaluations} & No submit button while demos are still running. Complete drafts can then be submitted together. \textbf{With an instructor:} goes to \code{SUBMITTED} for review. \textbf{Without one:} \ui{Release marks} asks for confirmation, then goes straight to \code{FINALIZED} and the student sees it. Toasts say ``released'' or ``sent for review''.\\
  T11 & Fixes a returned evaluation\newline{\color{muted}\footnotesize Marking sheet} & \code{saveDraft}, \code{submitEvaluations} & \code{evaluation.returned} notification includes the instructor's comment, which the sheet shows above the scores.\\
  T12 & Closes / archives\newline{\color{muted}\footnotesize \ui{Booking} panel, \code{…/manage/settings}} & \code{closeAssignment}, \code{setCourseArchived} & Close stops new bookings and changes, but existing bookings stay. Archive hides the course from dashboards and makes it read-only until restored. Settings and archiving belong to the instructor, or to a TA when the course has none.\\
  \bottomrule
\end{xltabular}}

\section{Instructor}

Everything a TA can do, plus:

{\tablefont
\begin{xltabular}{\linewidth}{@{}l P{34mm} P{30mm} L@{}}
  \toprule
  \thead{\#} & \thead{Step and screen} & \thead{Service} & \thead{What happens / feedback}\\
  \midrule
  \endhead
  I1 & Watches progress\newline{\color{muted}\footnotesize \code{…/manage}} & \code{reports.courseProgress} & One card per assignment: its booking state, the next thing to do in words (e.g. ``3 past demos need attendance'', ``5 demos to mark''), and two labelled bars, \ui{Booked} and \ui{Marked}, each with ``x of y''.\\
  I2 & Gives the final check\newline{\color{muted}\footnotesize \ui{Final check} task on the home page, \code{…/manage/review}} & \code{reviewEvaluation}, \code{finalizeMany} & \textbf{Finalize} releases marks to the student, whose email lists the total, every rubric row and the feedback. \textbf{Return} needs a comment, which goes to the TA. Marks an instructor submits themselves are final straight away.\\
  I3 & Corrects a released mark\newline{\color{muted}\footnotesize Student page} & \code{unlockEvaluation}, \code{saveDraft}, \code{submitEvaluations} & \ui{Unlock to edit} needs a reason. The marks reopen on the same page; \ui{Release corrected marks} finalizes them again, and the student gets \ui{Marks updated} with the new total and the old one.\\
  I4 & Exports\newline{\color{muted}\footnotesize \ui{Students \& marks} tab → \ui{Export CSV}} & \code{reports.exportAssignmentCsv} & One row per student: slot, TA, venue, attendance, rubric scores, total, status. Formula-injection safe.\\
  \bottomrule
\end{xltabular}}

\section{Admin}

Admins run the platform from the web console. They are never students or staff of a course (enrolling an admin is refused everywhere), don't book anything, and don't need offline access. They always land on \route{/admin}, and the header shows only \ui{Administration}. Their account page has just name, clock, password and devices.

{\tablefont
\begin{xltabular}{\linewidth}{@{}l P{34mm} P{30mm} L@{}}
  \toprule
  \thead{\#} & \thead{Step and screen} & \thead{Service} & \thead{What happens / feedback}\\
  \midrule
  \endhead
  A0 & Checks what needs attention\newline{\color{muted}\footnotesize \route{/admin} (Overview)} & \code{adminOverview} & \ui{Needs attention} lists only real problems, each with a link to fix it: failed emails, emails stuck in the queue for more than 15 minutes (the worker is down), active courses with no instructor or TA, and invites not accepted after 7 days. Below that, platform numbers (users, activation, courses, demos in the next 7 days, email) and recent activity.\\
  A1 & Finds or invites a person\newline{\color{muted}\footnotesize \route{/admin/users}} & \code{adminListUsers}, \code{adminInviteUser} & \ui{Invite someone}: email, name, and optionally admin; they get an invite and can then create a course or be added to one. Search plus filters: active, awaiting activation, invite older than 7 days, disabled, admins, not in any course.\\
  A1b & Helps a person\newline{\color{muted}\footnotesize \route{/admin/users/[id]}} & \code{adminGetUser}, \code{adminSendAccessEmail}, \code{adminUpdateUser}, \code{adminSetUserDisabled}, \code{adminSetAdmin} & Account status, last sign-in or invite age, courses and roles, upcoming demos, history (done to them and by them). \ui{Resend invite} or \ui{Send password reset}. Fix the name, or a mistyped email while the account is still invited (a new invite goes to the corrected address). Disable (ends sessions and can release upcoming bookings) or enable. Make or remove admin; someone who is in a course can't be made admin until they're removed from it. Admins can't disable themselves or remove their own admin rights.\\
  A2 & Keeps courses staffed\newline{\color{muted}\footnotesize \route{/admin/courses}} & \code{adminCreateCourse}, \code{adminListCourses}, \code{adminAssignStaff}, \code{adminSetCourseArchived} & \ui{New course}: code, title, term and the instructor or TA who runs it (invited if new). Filters: active, no staff, archived. Each course shows its staff, students, open assignments and upcoming demos. \ui{Assign staff} adds an instructor or TA by email (never an admin): new people get an invite, existing ones a notification. Archive or restore. Opening a course shows its manage pages with full rights (``Viewing as admin'').\\
  A3 & Fixes email delivery\newline{\color{muted}\footnotesize \route{/admin/email} (badge = failed count)} & \code{adminEmailQueue}, \code{adminRetryEmails} & Failed, waiting and sent emails with the SMTP error. Retry one, or all after fixing the cause.\\
  A4 & Audits\newline{\color{muted}\footnotesize \route{/admin/audit}} & \code{adminAuditLog} & Readable sentences (``Ada Admin assigned staff to MATH 210'') with links to users and courses, filtered by area or person, 100 per page, with before/after details on demand.\\
  \bottomrule
\end{xltabular}}


\section{Platform team}\label{sec:journey-platform}

The product team (whoever builds Slotty) reads the evidence for what to improve. That is not the
admins' job: admins run the deployment. A platform account is a \textbf{dedicated account}: never
an admin, never in a course, granted only from the command line (\code{bun run platform:grant},
\cref{sec:ops-platform}), and invisible to the admin console (ADR-016). It lands on
\route{/platform}; the header shows only \ui{Platform}. It isn't tracked, surveyed or offered the
feedback button.

{\tablefont
\begin{xltabular}{\linewidth}{@{}l P{34mm} P{30mm} L@{}}
  \toprule
  \thead{\#} & \thead{Step and screen} & \thead{Service} & \thead{What happens / feedback}\\
  \midrule
  \endhead
  P1 & Looks for evidence of what to improve\newline{\color{muted}\footnotesize \route{/platform/insights}} & \code{usageOverview}, \code{workflowMetrics}, \code{funnels}, \code{performance}, \code{clientErrors}, \code{surveys} & Six reports over 7, 30 or 90 days. \ui{Usage}: people active this week by role, people per day, pages by views (the bottom of the list is what hardly anyone opens), features used, and a 13-month trend. \ui{Booking \& marking}: reschedules, cancellations (and how many were rebooked within 30 minutes), no-shows, waitlist conversion, requests and how fast they were answered, demo-to-marks time, unlocks; these come from the booking tables, so they count everyone. \ui{Funnels}: visits that saw the slot picker or marking sheet, pressed submit, succeeded or got an error, by device. \ui{Speed}: Core Web Vitals at the 75th percentile by device, and the slowest pages. \ui{Errors}: browser errors grouped by message, with pages and releases. \ui{Surveys}: answers to the one-question surveys. See \cref{sec:rules-usage}.\\
  P2 & Reads feedback\newline{\color{muted}\footnotesize \route{/platform/feedback} (badge = new count)} & \code{listFeedback}, \code{setFeedbackStatus} & Anonymous messages from the feedback button with rating, role, page type and release. \ui{Mark triaged} once it's noted somewhere (a roadmap item), \ui{Mark done} once acted on, \ui{Reopen}. Status changes are audited under \code{platform.*}, which the admin audit log leaves out.\\
  \bottomrule
\end{xltabular}}

\section{Feedback and surveys (students and staff)}\label{sec:journey-feedback}

Everyone except admins has a \ui{Send feedback} button (speech-bubble icon) in the header: a
message and an optional 1--5 rating, sent anonymously (\code{submitFeedback}). After a student
books or moves a demo, and after staff submit marks, a small card in the corner asks one question
(``How easy was it to book your demo?'') with answers 1--5 or \ui{Not now}; each question is asked
at most once every 14 days (\code{claimSurvey}, \code{answerSurvey}). \route{/account} has
\ui{Usage data}, on by default, and \route{/privacy} explains what is collected
(\cref{sec:rules-usage}).

\section{End-to-end lifecycle (one assignment)}

\begin{figure}[htbp]
  \centering
  % One assignment's lifecycle (chapter 3): four lanes, read left to right, top to bottom.
% Laid out on a grid (column × 27 mm, lane × 17 mm) so no arrow crosses another.
\begin{tikzpicture}[x=27mm, y=-17mm,
    lcstep/.style={svc=21mm, minimum height=9mm, font=\sffamily\scriptsize},
    lane/.style={font=\sffamily\scriptsize\bfseries, text=muted, anchor=east, align=right}]
  % lane 0: staff set up
  \node[lcstep] (course)  at (0,0) {Create course};
  \node[lcstep] (roster)  at (1,0) {Import roster};
  \node[lcstep] (assign)  at (2,0) {Assignment, policy, rubric};
  \node[lcstep] (slots)   at (3,0) {Add slots (draft)};
  \node[lcstep] (publish) at (4,0) {Open booking};
  % lane 1: students
  \node[lcstep] (book)    at (2,1) {Book};
  \node[lcstep] (change)  at (3,1) {Reschedule or cancel};
  \node[lcstep] (remind)  at (4,1) {Reminders 24 h and 1 h};
  % lane 2: demo day
  \node[lcstep] (today)   at (2,2) {Demo day: now and next};
  \node[lcstep] (mark)    at (3,2) {Attendance and marks};
  \node[lcstep] (submit)  at (4,2) {Submit};
  % lane 3: results
  \node[lcstep] (review)  at (3,3) {Instructor review};
  \node[lcstep] (fin)     at (4,3) {Finalized};
  \node[lcstep] (see)     at (2,3) {Student sees marks};
  \node[lcstep] (close)   at (1,3) {Export, close, archive};

  \node[lane] at (-.5,0) {Staff};
  \node[lane] at (-.5,1) {Students};
  \node[lane] at (-.5,2) {Demo day};
  \node[lane] at (-.5,3) {Results};

  \draw[flow] (course) -- (roster);
  \draw[flow] (roster) -- (assign);
  \draw[flow] (assign) -- (slots);
  \draw[flow] (slots) -- (publish);
  \draw[async] (publish.south) -- ++(0,.22) -| node[edgelabel, pos=.25]{students notified} (book.north);
  \draw[flow] (book) -- (change);
  \draw[async] (change) -- (remind);
  \draw[flow] (remind.south) -- ++(0,.22) -| (today.north);
  \draw[flow] (today) -- (mark);
  \draw[flow] (mark) -- (submit);
  \draw[flow] (submit) -- node[edgelabel, right]{no instructor} (fin);
  \draw[flow] (submit.south west) -- node[edgelabel, sloped, below]{to review} (review.north east);
  \draw[flow] (review) -- node[edgelabel, below]{finalize} (fin);
  \draw[ctrl] (review) -- node[edgelabel, left]{return} (mark);
  \draw[flow] (fin.south) -- ++(0,.25) -| (see.south);
  \draw[flow] (see) -- (close);
\end{tikzpicture}

  \caption{One assignment from setup to archive.}
  \label{fig:assignment-lifecycle}
\end{figure}


Planned changes to these flows are listed in \cref{ch:roadmap}.
